\chapter{Surapprentissage}\label{ch:rs:overfitting}

Mille réglages d'un système de tendance à moyennes mobiles sont soumis à un \omterm{def:rs:vectorised-backtests:backtest}{backtest} sur dix-huit ans de rendements quotidiens. Le meilleur affiche un ratio de Sharpe
de 0.44~; sur les cinq années suivantes, il perd, avec un ratio de Sharpe de $-0.32$. Les rendements n'étaient que du bruit pur, et les mille systèmes constituent une recherche menée à travers
celui-ci. On lance la même recherche sur des rendements comportant une tendance réelle, injectée à dessein, et le meilleur système de la période de recherche déçoit de nouveau par la suite, rapportant moins
que le système moyen~: la tendance était présente pour tous, et la recherche a sélectionné du bruit venant s'y ajouter. Le surapprentissage en finance tient rarement à un modèle
à trop de paramètres~; il s'agit plutôt d'un choix parmi trop de \omterm{def:rs:vectorised-backtests:backtest}{backtests}, mesuré sur un historique trop court. Ce chapitre le mesure (la \omterm{def:rs:overfitting:pbo}{probabilité de
surapprentissage du backtest}), le borne (la \omterm{def:rs:overfitting:minbtl}{longueur minimale du backtest}) et prévient la fuite qui le rend invisible dans la validation croisée (\omterm{def:rs:overfitting:purging}{purge} et
\omterm{def:rs:overfitting:purging}{embargo}), avec \texttt{firm.overfit}.

\section{Sélection parmi de nombreux essais}

\begin{definition}[Surapprentissage du backtest, dans l'échantillon, hors échantillon]\label{def:rs:overfitting:overfitting}
Les données \emph{dans l'échantillon}\index{dans l'échantillon} (in-sample) sont celles utilisées pour construire ou choisir une stratégie~; les données \emph{hors échantillon}\index{hors échantillon} (out-of-sample) sont celles
qu'elle n'a pas vues. \emph{surapprentissage du backtest}\index{surapprentissage du backtest} (\omterm{def:rs:vectorised-backtests:backtest}{backtest} overfitting) est la sélection d'une stratégie dont la performance dans l'échantillon tient davantage à
l'échantillon particulier qu'à une propriété qui persiste, de sorte que sa performance hors échantillon reste en deçà de celle qui avait motivé la sélection.
\end{definition}

Les systèmes du chapitre prennent une position acheteuse sur le marché lorsqu'une moyenne mobile rapide du cours est au-dessus d'une moyenne lente, et vendeuse sinon~: fenêtres rapides de 2 à 40
jours, fenêtres lentes de 50 à 540, soit mille couples. Sur 25 ans de rendements quotidiens synthétiques, les 20 premières années servent à la recherche (17.9 ans une fois que la
moyenne la plus longue a démarré), les cinq dernières servent de test. Sur du pur bruit, le ratio de Sharpe \omterm{def:rs:overfitting:overfitting}{dans l'échantillon} du meilleur système est de 0.44, son ratio \omterm{def:rs:overfitting:overfitting}{hors échantillon} est de $-0.32$~; pour le système médian
le ratio \omterm{def:rs:overfitting:overfitting}{dans l'échantillon} vaut 0.17. Le meilleur attendu parmi mille stratégies nulles indépendantes sur 17.9 ans serait de 0.77 (livre~4, chapitre~12)~; les systèmes
sont loin d'être indépendants (leurs ratios de Sharpe ont un écart-type de seulement 0.076 autour d'un niveau commun), et le maximum attendu issu
de cette dispersion dépasse le niveau commun de 0.25, soit à peu près ce qui a été observé.

Avec une tendance injectée, une dérive persistante du genre de celles que gagnent les suiveurs de tendance (une demi-vie de 120 jours), chaque système la capte~: le ratio de Sharpe médian \omterm{def:rs:overfitting:overfitting}{dans l'échantillon}
est de 0.50, et la moyenne \omterm{def:rs:overfitting:overfitting}{hors échantillon} de 0.26. Le meilleur système \omterm{def:rs:overfitting:overfitting}{dans l'échantillon}, à 0.77, rapporte 0.14 \omterm{def:rs:overfitting:overfitting}{hors échantillon}, moins que le système moyen. Le classement n'est pas inutile (entre les systèmes, les ratios de Sharpe \omterm{def:rs:overfitting:overfitting}{dans l'échantillon} et
\omterm{def:rs:overfitting:overfitting}{hors échantillon} sont corrélés à 0.33, et les cinquante meilleurs \omterm{def:rs:overfitting:overfitting}{dans l'échantillon} rapportent 0.29 \omterm{def:rs:overfitting:overfitting}{hors échantillon}), mais c'est sur le gagnant unique que se concentre le bruit~: la
recherche a trouvé la tendance, et sélectionner un seul point de celle-ci a conduit à retenir du bruit. Sur du bruit, la corrélation vaut $-0.05$ et les cinquante meilleurs rapportent $-0.02$.

\section{La probabilité de surapprentissage du backtest}

\begin{definition}[Probabilité de surapprentissage du backtest, validation croisée combinatoire symétrique]\label{def:rs:overfitting:pbo}
La \emph{probabilité de surapprentissage du backtest}\index{probabilité de surapprentissage du backtest} (PBO, probability of \omterm{def:rs:vectorised-backtests:backtest}{backtest} overfitting) d'une procédure de sélection est la probabilité que la
stratégie qu'elle sélectionne comme meilleure \omterm{def:rs:overfitting:overfitting}{dans l'échantillon} se situe sous la médiane des candidats \omterm{def:rs:overfitting:overfitting}{hors échantillon}. \emph{Validation croisée
combinatoire symétrique}\index{validation croisée combinatoire symétrique} (combinatorially symmetric cross-validation, CSCV) l'estime en découpant l'historique en $S$ blocs et, pour chaque choix de
$S/2$ blocs comme ensemble \omterm{def:rs:overfitting:overfitting}{dans l'échantillon}, en classant le gagnant \omterm{def:rs:overfitting:overfitting}{dans l'échantillon} parmi tous les candidats sur l'autre moitié~; le PBO est la part des découpages pour lesquels son rang
est inférieur à la médiane (Bailey, Borwein, López de Prado et Zhu).
\end{definition}

Le rang $\omega$ du gagnant parmi $N$ candidats, exprimé sous forme de fraction, est transformé en logit $\lambda = \ln(\omega/(1 - \omega))$~: positif lorsque le gagnant
reste au-dessus de la médiane, négatif lorsqu'il passe en dessous. Avec 16 blocs, il existe 12\,870 découpages~; \texttt{firm.overfit.cscv} en utilise 3\,000, tirés au
hasard, et calcule les ratios de Sharpe de chaque découpage comme produits matriciels de sommes par blocs. Le PBO de la recherche à mille systèmes vaut 0.82 sur du bruit et 0.74 avec la
tendance injectée (\cref{fig:rs:overfitting:logits})~; la régression des ratios de Sharpe \omterm{def:rs:overfitting:overfitting}{hors échantillon} des gagnants sur ceux \omterm{def:rs:overfitting:overfitting}{dans l'échantillon} présente des pentes de $-0.89$
et $-0.91$. Un PBO supérieur à un demi signifie que la sélection fait moins bien qu'un choix au hasard~: parmi des systèmes corrélés, le gagnant \omterm{def:rs:overfitting:overfitting}{dans l'échantillon} est celui qui a le mieux ajusté
le bruit de sa moitié, et le bruit de l'autre moitié en est le miroir.

\begin{omfigure}
\begin{tikzpicture}
\begin{axis}[width=11cm,height=5.4cm,ybar,bar width=7pt,xlabel={logit du rang hors échantillon du gagnant dans l'échantillon},ylabel={part des découpages},
  tick label style={font=\small},label style={font=\small},xmin=-6.2,xmax=6.2,ymin=0,grid=major,grid style={gray!25},
  yticklabel style={/pgf/number format/fixed,/pgf/number format/precision=2},scaled y ticks=false,legend columns=2,
  legend style={font=\footnotesize,at={(0.5,-0.3)},anchor=north,draw=none,fill=none,
  /tikz/every even column/.append style={column sep=8pt}},table/col sep=comma]
\addplot[fill=gray!45,draw=gray] table[x=x,y=share]{figdata/research/20-overfitting/logits_noise.csv};
\addplot[fill=omDef!55,draw=omDef] table[x=x,y=share]{figdata/research/20-overfitting/logits_trend.csv};
\draw[omThm,thick,dashed] (axis cs:0,0) -- (axis cs:0,0.2);
\legend{bruit (PBO 0.82),tendance injectée (PBO 0.74)}
\end{axis}
\end{tikzpicture}

\omcaption{CSCV sur les années de recherche des mille systèmes de tendance (16 blocs, 3\,000 découpages)~: la distribution du logit du
rang \omterm{def:rs:overfitting:overfitting}{hors échantillon} du gagnant \omterm{def:rs:overfitting:overfitting}{dans l'échantillon}. La masse située à gauche de zéro est la \omterm{def:rs:overfitting:pbo}{probabilité de surapprentissage du backtest}. Données~: \texttt{rs\_overfit} sur des rendements synthétiques.}\label{fig:rs:overfitting:logits}
\end{omfigure}

\begin{omfigure}
\begin{tikzpicture}
\begin{axis}[name=a,width=6.6cm,height=5.6cm,title={\small bruit},xlabel={\small ratio de Sharpe dans l'échantillon},ylabel={\small ratio de Sharpe hors échantillon},
  tick label style={font=\footnotesize},xmin=-0.3,xmax=0.9,ymin=-0.8,ymax=1.2,grid=major,grid style={gray!25},table/col sep=comma]
\addplot[only marks,mark=*,mark size=0.8pt,gray] table[x=is,y=oos]{figdata/research/20-overfitting/scatter_noise.csv};
\end{axis}
\begin{axis}[at={(a.east)},anchor=west,xshift=1.2cm,width=6.6cm,height=5.6cm,title={\small tendance injectée},xlabel={\small ratio de Sharpe dans l'échantillon},
  tick label style={font=\footnotesize},xmin=-0.3,xmax=0.9,ymin=-0.8,ymax=1.2,grid=major,grid style={gray!25},table/col sep=comma]
\addplot[only marks,mark=*,mark size=0.8pt,omDef] table[x=is,y=oos]{figdata/research/20-overfitting/scatter_trend.csv};
\end{axis}
\end{tikzpicture}

\omcaption{Ratios de Sharpe \omterm{def:rs:overfitting:overfitting}{dans l'échantillon} (17.9 ans) contre \omterm{def:rs:overfitting:overfitting}{hors échantillon} (5 ans) des systèmes de tendance (un sur quatre parmi les mille). Sur du bruit, les nuages de points sont sans lien (corrélation $-0.05$)~; avec la tendance, ils sont corrélés à 0.33, pourtant le système en tête \omterm{def:rs:overfitting:overfitting}{dans l'échantillon} ne figure pas parmi les meilleurs \omterm{def:rs:overfitting:overfitting}{hors échantillon}. Données~: \texttt{rs\_overfit.systems}.}\label{fig:rs:overfitting:scatter}
\end{omfigure}

\section{Longueur minimale du backtest}

\begin{definition}[Longueur minimale du backtest]\label{def:rs:overfitting:minbtl}
La \emph{longueur minimale du backtest}\index{longueur minimale du backtest} (minimum \omterm{def:rs:vectorised-backtests:backtest}{backtest} length) pour $N$ essais et un ratio de Sharpe cible est le nombre d'années d'historique en deçà duquel le
meilleur attendu de $N$ stratégies sans habileté affiche \omterm{def:rs:overfitting:overfitting}{dans l'échantillon} un ratio de Sharpe au moins égal à la cible.
\end{definition}

\begin{proposition}[Quel historique une recherche exige]\label{prop:rs:overfitting:minbtl}
Si $N$ stratégies indépendantes sans habileté sont chacune mesurées sur $y$ années, leurs estimations annuelles du ratio de Sharpe ont une variance d'environ $1/y$, et le
maximum attendu croît comme $\sqrt{2\ln N/y}$~; un ratio de Sharpe cible $s$ est censé être atteint par hasard, sauf si $y \gtrsim 2\ln N/s^2$.
\end{proposition}

\begin{proof}
Le maximum attendu de $N$ normales centrées réduites croît comme $\sqrt{2\ln N}$ (livre~4, chapitre~12)~; on passe à l'échelle de l'erreur type $1/\sqrt y$ et l'on résout
$\sqrt{2\ln N/y} = s$.
\end{proof}

Avec le maximum attendu calculé exactement (\texttt{firm.multitest.expected\_max\_sr}), mille essais nécessitent 10.6 ans avant qu'un ratio de Sharpe de 1 cesse
d'être imputable à la chance, 42.4 ans pour 0.5, et 3.3 ans pour 1.8 (\cref{fig:rs:overfitting:minbtl}). La borne est conservatrice lorsque les essais sont
corrélés, comme le sont les systèmes de tendance, et elle sert de repère lorsqu'ils ne le sont pas~; son enseignement porte sur l'échelle~: quatre fois plus d'historique pour un ratio de Sharpe deux fois moindre. Harvey et
Liu (2015) a transformé la même logique en décote applicable à tout ratio de Sharpe publié, compte tenu du nombre de tests qui le sous-tendent.

\begin{omfigure}
\begin{tikzpicture}
\begin{semilogxaxis}[width=10cm,height=5.6cm,xlabel={nombre d'essais},ylabel={années d'historique},tick label style={font=\small},label style={font=\small},
  xmin=2,xmax=10000,ymin=0,ymax=60,grid=major,grid style={gray!25},legend pos=north west,legend style={font=\footnotesize},table/col sep=comma]
\addplot[omThm,very thick] table[x=trials,y=sr05]{figdata/research/20-overfitting/minbtl.csv};
\addplot[omDef,very thick] table[x=trials,y=sr10]{figdata/research/20-overfitting/minbtl.csv};
\addplot[black!60,thick,dashed] table[x=trials,y=sr18]{figdata/research/20-overfitting/minbtl.csv};
\legend{Sharpe cible 0.5,Sharpe cible 1.0,Sharpe cible 1.8}
\end{semilogxaxis}
\end{tikzpicture}

\omcaption{\omterm{def:rs:overfitting:minbtl}{Longueur minimale du backtest}~: le nombre d'années d'historique en deçà duquel le meilleur résultat attendu de $N$ stratégies nulles indépendantes atteint le ratio de Sharpe cible.
Données~: \texttt{firm.overfit.min\_backtest\_length}.}\label{fig:rs:overfitting:minbtl}
\end{omfigure}

\section{Validation croisée purgée et avec embargo}

\begin{definition}[Chevauchement des étiquettes, purge, embargo, validation croisée combinatoire purgée]\label{def:rs:overfitting:purging}
Il y a \emph{chevauchement des étiquettes}\index{chevauchement des étiquettes} (label overlap) lorsque les cibles de différentes observations partagent des périodes, comme le font les rendements à terme de $h$ jours observés chaque jour.
La \emph{purge}\index{purge} (purging) retire de l'ensemble d'entraînement chaque observation dont l'étiquette chevauche dans le temps celles de l'ensemble de test~; un
\emph{embargo}\index{embargo} retire également les observations d'entraînement qui suivent immédiatement l'ensemble de test, car les variables explicatives sont corrélées en série.
La \emph{validation croisée combinatoire purgée}\index{validation croisée combinatoire purgée} (combinatorial purged cross-validation, CPCV) utilise chaque choix de $k$ plis de test parmi $n$, purgés et
avec embargo, en recombinant les plis de test en $\binom{n}{k}k/n$ chemins de \omterm{def:rs:vectorised-backtests:backtest}{backtest} (López de Prado, 2018~; telle qu'elle est implémentée dans des bibliothèques open source comme skfolio).
\end{definition}

La validation croisée $k$-fold avec mélange, celle que proposent par défaut les bibliothèques généralistes d'apprentissage automatique, place les jours voisins aussi bien dans les ensembles d'entraînement que dans ceux de test. Quand
les étiquettes se chevauchent, un modèle peut prédire l'étiquette d'un jour de test à partir des jours d'entraînement dont les étiquettes partagent ses périodes. La démonstration du chapitre est extrême à
dessein~: rendements à terme de 20 jours d'un pur bruit, prédits par les cinq plus proches voisins, sur une variable explicative qui évolue aussi lentement que le temps (une variable de régime), cinq
plis~:

\begin{center}
\small
\begin{tabular}{@{}lcccc@{}}
\toprule
validation croisée & mélangée & contiguë & purgée & purgée et avec \omterm{def:rs:overfitting:purging}{embargo}\\
\midrule
$R^2$ d'une prévision du bruit & 0.90 & $-0.63$ & $-0.57$ & $-0.93$\\
\bottomrule
\end{tabular}
\end{center}

Les plis mélangés indiquent que le modèle explique 90~\% de la variance du bruit~; tout schéma qui respecte le temps affirme qu'il fait moins bien qu'une constante. Les plis contigus
ne fuient qu'à leurs bords, la \omterm{def:rs:overfitting:purging}{purge} referme les bords, et l'\omterm{def:rs:overfitting:purging}{embargo} protège contre les variables explicatives qui transmettent de l'information d'un pli à l'autre. Avec un
schéma combinatoire, le même modèle peut être jugé sur de nombreux chemins plutôt que sur un seul, ce qui transforme un unique ratio de Sharpe \omterm{def:rs:overfitting:overfitting}{hors échantillon} en distribution.

\section{Analyse glissante}

\begin{definition}[Analyse glissante]\label{def:rs:overfitting:walkforward}
L'\emph{analyse glissante}\index{analyse glissante} (walk-forward analysis) répète dans le temps le cycle qui consiste à choisir une stratégie sur une fenêtre passée et à la trader sur la
période suivante, et juge les périodes \omterm{def:rs:overfitting:overfitting}{hors échantillon} recollées.
\end{definition}

L'\omterm{def:rs:overfitting:walkforward}{analyse glissante} est la procédure que suivrait réellement la société, et elle mesure la procédure plutôt qu'une stratégie. Refaire chaque année le choix du meilleur des mille
systèmes sur les cinq années précédentes donne un ratio de Sharpe recollé de $-0.04$ sur le bruit et de 0.29 avec la tendance implantée~: le ratio de Sharpe réel de la tendance,
à peu près ce que rapporte le système moyen, et deux fois ce qu'a rapporté le gagnant \omterm{def:rs:overfitting:overfitting}{dans l'échantillon} (in-sample) sur toute la recherche. Sa limite tient au chemin unique~: une seule histoire, une seule suite de
choix~; CPCV et CSCV existent parce qu'un seul chemin est un petit échantillon.

\section{Garde-fous pratiques}

Les garde-fous relèvent de l'organisation avant d'être statistiques. Compter chaque essai dans le \omterm{def:rs:the-research-process:log}{journal de recherche} (chapitre~1), y compris ceux qui ont été abandonnés~; reporter la PBO de la
recherche et le ratio de Sharpe déflaté de son gagnant (livre~4, chapitre~12)~; fixer le protocole d'évaluation avant de voir les résultats~; garder un \omterm{def:rs:the-research-process:prereg}{échantillon réservé} que personne ne touche
jusqu'à la décision~; privilégier la large zone de paramètres qui fonctionne plutôt que son meilleur point (un plateau, non un pic~: ici, le système moyen a battu le gagnant)~; et
exiger un mécanisme, car un mécanisme résiste au changement d'échantillon, contrairement à un chiffre ajusté. Arnott, Harvey et Markowitz (2019) exposent ces principes dans un
protocole de recherche en apprentissage automatique~; la PBO et la \omterm{def:rs:overfitting:minbtl}{longueur minimale du backtest} les chiffrent.

\section{Tutoriel~: mille systèmes de tendance}\label{tut:rs:overfitting:tut}

\begin{tutorial}
\textbf{Objectif.} Rechercher parmi mille systèmes de tendance sur du bruit et sur une tendance implantée, estimer la PBO par CSCV, procéder à une \omterm{def:rs:overfitting:walkforward}{analyse glissante}, calculer la \omterm{def:rs:overfitting:minbtl}{longueur minimale du backtest}, et
comparer quatre schémas de validation croisée sur des étiquettes qui se chevauchent. \textbf{État final~:}
Les \cref{fig:rs:overfitting:logits,fig:rs:overfitting:scatter,fig:rs:overfitting:minbtl}~; le tableau.
\begin{tutsteps}
\item \textbf{CSCV}~: sommes par blocs, puis ratios de Sharpe de chaque découpage calculés sous forme de produits matriciels, puis rang \omterm{def:rs:overfitting:overfitting}{hors échantillon} du gagnant.
\omcode{code/firm/overfit/firm_overfit.py}{33}{64}{Probabilité de surapprentissage du backtest par CSCV.}\label{lst:rs:overfitting:cscv}
\item \textbf{k-fold purgé avec \omterm{def:rs:overfitting:purging}{embargo}.}
\omcode{code/firm/overfit/firm_overfit.py}{67}{80}{Purge et embargo.}\label{lst:rs:overfitting:purge}
\item \textbf{Exécuter} \texttt{search()} pour les deux mondes, \texttt{walk\_forward\_sr()}, \texttt{minbtl()}, \texttt{label\_leak()} et \texttt{fig\_overfit.py}.
\end{tutsteps}
\textbf{Pour aller plus loin.} Choisir le centre du meilleur plateau de paramètres plutôt que son meilleur point isolé et comparer son ratio de Sharpe \omterm{def:rs:overfitting:overfitting}{hors échantillon}~; exécuter
CSCV sur cent stratégies indépendantes plutôt que sur mille stratégies corrélées, et observer la PBO se rapprocher d'un demi sur le bruit.
\end{tutorial}

\section{Construction~: la boîte à outils de surapprentissage}\label{bld:rs:overfitting:overfit}

\begin{build}
\bfield{Rôle} Chaque recherche menée par la société est présentée avec sa probabilité de surapprentissage, sa \omterm{def:rs:overfitting:minbtl}{longueur minimale du backtest} et une validation croisée qui respecte
l'ordre temporel.
\bfield{Interface} \texttt{cscv(M, n\_blocks, max\_splits, seed)}, \texttt{purged\_kfold(start, end, k, embargo)}, \texttt{cpcv(start, end, n\_folds, n\_test, embargo)},
\texttt{cpcv\_paths}, \texttt{walk\_forward(n, train, test, step, anchored)}, \texttt{min\_backtest\_length(n\_trials, sr\_target)}.
\bfield{Règles} Aucune validation croisée mélangée sur des séries temporelles~; les étiquettes portent leur date de début et de fin~; le \omterm{def:rs:the-research-process:log}{nombre d'essais} provient du \omterm{def:rs:the-research-process:log}{journal de recherche}, et non de la mémoire.
\bfield{Tests d'acceptation} \texttt{code/firm/overfit/tests/}~: une PBO proche d'un demi et une pente non positive sur le bruit, proche de zéro avec une stratégie performante~; \omterm{def:rs:overfitting:purging}{purge} et
\omterm{def:rs:overfitting:purging}{embargo} à la main~; les quinze découpages et les cinq chemins de la CPCV, ensembles d'entraînement et de test disjoints~; fenêtres de l'\omterm{def:rs:overfitting:walkforward}{analyse glissante}, ancrées et roulantes~; la monotonie de la \omterm{def:rs:overfitting:minbtl}{longueur minimale
du backtest} et son quadruplement pour une cible divisée par deux.
\bfield{Pour aller plus loin} Distributions des ratios de Sharpe en CPCV pour un modèle~; ratio de Sharpe déflaté à partir du nombre effectif d'essais~; tableau de bord de la PBO alimenté par le \omterm{def:rs:the-research-process:log}{journal de recherche}.
\end{build}

\begin{omsources}
\item D.~H.~Bailey, J.~M.~Borwein, M.~López de Prado et Q.~J.~Zhu, ``The probability of backtest overfitting'', \emph{Journal of Computational Finance}, 2016.
\item C.~R.~Harvey et Y.~Liu, ``Backtesting'', \emph{Journal of Portfolio Management} 42(1), 2015.
\item M.~López de Prado, \emph{Advances in Financial Machine Learning}, Wiley, 2018~; skfolio, \texttt{CombinatorialPurgedCV} (documentation de la purge et de l'embargo).
\item R.~Arnott, C.~R.~Harvey et H.~Markowitz, ``A backtesting protocol in the era of machine learning'', \emph{Journal of Financial Data Science} 1(1), 2019.
\end{omsources}

\section{Exercices}

\begin{exercise}[$\star$]\label{exo:rs:overfitting:1}
Dans un découpage CSCV, le gagnant \omterm{def:rs:overfitting:overfitting}{dans l'échantillon} (in-sample) se classe 380e sur 1\,000 \omterm{def:rs:overfitting:overfitting}{hors échantillon} (1 étant le plus mauvais). Calculer son logit et dire si le découpage entre dans le calcul de la PBO.
\end{exercise}

\begin{exercise}[$\star$]\label{exo:rs:overfitting:2}
Calculer le nombre de chemins de \omterm{def:rs:vectorised-backtests:backtest}{backtest} que fournit la CPCV avec 10 plis dont 2 plis de test, puis avec 6 plis dont 3 plis de test.
\end{exercise}

\begin{exercise}[$\star$]\label{exo:rs:overfitting:3}
En utilisant l'approximation de la proposition, calculer le nombre d'années d'historique dont une recherche portant sur 100 essais a besoin pour un ratio de Sharpe cible de 1.
\end{exercise}

\begin{exercise}[$\star\star$]\label{exo:rs:overfitting:4}
Les étiquettes sont des rendements à terme de 10 jours observés chaque jour~; le pli de test couvre les jours 200 à 299. Déterminer quels jours d'entraînement la \omterm{def:rs:overfitting:purging}{purge} retire, et ceux qu'ajoute un \omterm{def:rs:overfitting:purging}{embargo} de
5 jours.
\end{exercise}

\begin{exercise}[$\star\star$]\label{exo:rs:overfitting:5}
Pourquoi la PBO peut-elle dépasser la moitié~? Expliquer à l'aide des mille systèmes de tendance corrélés.
\end{exercise}

\begin{exercise}[$\star\star$]\label{exo:rs:overfitting:6}
Pourquoi la tendance insérée n'a-t-elle pas fait du gagnant \omterm{def:rs:overfitting:overfitting}{dans l'échantillon} un bon choix, alors que chaque système a profité de la tendance~?
\end{exercise}

\begin{exercise}[$\star\star\star$]\label{exo:rs:overfitting:7}
\emph{Programmation.} Lancer la CSCV sur les 100 premiers des mille systèmes seulement, puis sur 100 stratégies de bruit indépendantes. Comparer les PBO obtenues et expliquer.
\end{exercise}

\begin{exercise}[$\star\star\star$]\label{exo:rs:overfitting:8}
\emph{Trouver la faille.} «~Nous avons réglé les hyperparamètres du modèle par validation croisée à 10 plis avec mélange, ce qui est la norme, de sorte que l'estimation \omterm{def:rs:overfitting:overfitting}{hors échantillon} est
sans biais.~»
\end{exercise}

\section{Problème~: Mille systèmes de tendance}

\begin{problem}[{Problème du week-end --- une recherche mesurée}]\label{pb:rs:overfitting:1}
Les mille systèmes de tendance du chapitre, appliqués à 25 ans de rendements synthétiques, sur du bruit et avec une tendance insérée.

\medskip\noindent\textbf{Partie I --- La recherche.}
\begin{enumerate}
\item Quel est le meilleur ratio de Sharpe \omterm{def:rs:overfitting:overfitting}{dans l'échantillon}, et quelle est sa valeur \omterm{def:rs:overfitting:overfitting}{hors échantillon} sur du bruit~? Et avec la tendance~?
\item Quelles sont les médianes des ratios de Sharpe \omterm{def:rs:overfitting:overfitting}{dans l'échantillon} et les moyennes \omterm{def:rs:overfitting:overfitting}{hors échantillon} dans chaque monde~?
\item Quel maximum donneraient mille stratégies nulles indépendantes sur 17.9~ans, et pourquoi le gagnant observé est-il inférieur~?
\item Que fait la tendance à chaque système, et qu'apporte la recherche~?
\end{enumerate}

\medskip\noindent\textbf{Partie II --- La PBO.}
\begin{enumerate}[resume]
\item Comment la CSCV calcule-t-elle rapidement les ratios de Sharpe d'une partition~?
\item Quelles PBO obtient-on, et quelles sont les pentes de dégradation dans chaque monde~?
\item Que signifie une PBO supérieure à un demi~?
\item Quelle serait la PBO pour une stratégie dotée d'un véritable edge parmi du bruit~?
\end{enumerate}

\medskip\noindent\textbf{Partie III --- Longueur et fuite d'information.}
\begin{enumerate}[resume]
\item Donner les longueurs minimales de \omterm{def:rs:vectorised-backtests:backtest}{backtest} pour mille essais, pour des cibles de 0.5, 1 et 1.8.
\item Pourquoi diviser la cible par deux quadruple-t-il la longueur~?
\item Quel $R^2$ donnent les quatre schémas de validation croisée pour la prévision du bruit~?
\item Pourquoi les plis mélangés fuient-ils, et que retirent respectivement la \omterm{def:rs:overfitting:purging}{purge} et l'\omterm{def:rs:overfitting:purging}{embargo}~?
\end{enumerate}

\medskip\noindent\textbf{Partie IV --- Le verdict.}
\begin{enumerate}[resume]
\item Quel ratio de Sharpe donne une resélection annuelle en walk-forward dans chaque monde~?
\item Énoncer le \emph{résultat nommé}~: la PBO de la recherche et l'espérance \omterm{def:rs:overfitting:overfitting}{hors échantillon} déflatée du système retenu.
\item Lequel faudrait-il trader~: le gagnant \omterm{def:rs:overfitting:overfitting}{dans l'échantillon}, le système moyen ou la procédure walk-forward~?
\item Que doit contenir le \omterm{def:rs:the-research-process:log}{journal de recherche} pour cette recherche~?
\item Comment un mécanisme modifierait-il la confiance qu'on accorde au résultat~?
\item Qu'apporte la décote de Harvey et Liu à la PBO~?
\item En quoi la CPCV améliore-t-elle l'\omterm{def:rs:overfitting:walkforward}{analyse glissante}~?
\item En une phrase~: qu'est-ce que le \omterm{def:rs:overfitting:overfitting}{surapprentissage du backtest} dans une recherche~?
\end{enumerate}
\end{problem}

\section{Questions d'entretien}

\begin{interviewq}[$\star$ \iqroles{chercheur}]\label{iq:rs:overfitting:1}
Vous avez testé 500 stratégies, et la meilleure affiche un ratio de Sharpe de 2 dans le \omterm{def:rs:vectorised-backtests:backtest}{backtest}. Que faites-vous ensuite~?
\end{interviewq}

\begin{interviewq}[$\star\star$ \iqroles{chercheur, ingénieur ML}]\label{iq:rs:overfitting:2}
Pourquoi la validation croisée k-fold classique est-elle inadaptée aux séries temporelles financières, et par quoi la remplacez-vous~?
\end{interviewq}

\begin{interviewq}[$\star\star$ \iqroles{chercheur}]\label{iq:rs:overfitting:3}
Expliquez la \omterm{def:rs:overfitting:pbo}{probabilité de surapprentissage du backtest} et la manière dont la CSCV l'estime.
\end{interviewq}

\begin{interviewq}[$\star\star$ \iqroles{chercheur}]\label{iq:rs:overfitting:4}
Quelle doit être la longueur d'un \omterm{def:rs:vectorised-backtests:backtest}{backtest} pour qu'on puisse se fier à un ratio de Sharpe de 1 trouvé parmi 1\,000 essais~?
\end{interviewq}

\begin{interviewq}[$\star\star$ \iqroles{chercheur, risque}]\label{iq:rs:overfitting:5}
Quelles pratiques d'organisation réduisent le surapprentissage dans une équipe de recherche~?
\end{interviewq}

\begin{interviewq}[$\star\star\star$ \iqroles{chercheur}]\label{iq:rs:overfitting:6}
Dérivez la loi d'échelle de la \omterm{def:rs:overfitting:minbtl}{longueur minimale du backtest} en fonction du \omterm{def:rs:the-research-process:log}{nombre d'essais} et du ratio de Sharpe cible.
\end{interviewq}
